%% kapa: exam=1 round=34 year=2023 period=2 form=A subject=accounting subject_ko=회계학 seq=7 number=47
%% source: Q-Net 감정평가사 기출문제 (cst00309) artlSeq=5211789 file=raw/1008/5211789/2023년도 제34회 감정평가사 1차시험 2교시.pdf
%% conversion: kordoc 4.24.51 -> tools/gen1.py (auto); review=proofread-sonnet; flags=none
\begin{question}{47}
재무정보의 질적 특성에 관한 설명으로 옳지 않은 것을 모두 고른 것은?
\begin{framedtext}
ㄱ. 오류가 없다는 것은 현상의 기술에 오류나 누락이 없고, 보고 정보를 생산하는 데 사용되는 절차의 선택과 적용 시 절차 상 완벽하게 정확하다는 것을 의미한다. \par
ㄴ. 재무정보가 과거 평가에 대해 피드백을 제공한다면 확인가치를 갖는다. \par
ㄷ. 회계기준위원회는 중요성에 대한 획일적인 계량 임계치를 정하거나 특정한 상황에서 무엇이 중요한 것인지를 미리 결정할 수 있다. \par
ㄹ. 목적적합하고 충실하게 표현된 정보의 유용성을 보강시키는 질적 특성으로는 비교가능성, 검증가능성, 적시성 및 이해가능성이 있다. \par
\end{framedtext}
\begin{choices}
  \item ㄱ, ㄴ
  \item ㄱ, ㄷ
  \item ㄱ, ㄹ
  \item ㄴ, ㄷ
  \item ㄷ, ㄹ
\end{choices}
\end{question}
\answer{2}
